\documentclass[11pt,a4paper]{article}
\usepackage{auditstyle}
\lhead{Fresh Glimm--Jaffe proof audit}
\rhead{16 August 2026}
\usepackage{bm}

\newcommand{\eps}{\varepsilon}
\newcommand{\W}{\mathcal W}
\newcommand{\T}{\mathbb T}
\newcommand{\dd}{\,d}
\DeclareMathOperator{\gap}{gap}

\title{Part II finite-scale model sheet\\[3pt]
\large MMST nearest-neighbour kinematics, the Daubechies \(D4\) scaling map, and a fixed Wick-quartic convention}
\author{Lamport-form definition and consistency proof}
\date{16 August 2026}

\begin{document}
\maketitle

\begin{resultbox}
\textbf{Scope fixed.}  The resolution limit is first formulated on one fixed spatial
torus \(\T_L=\R/(2L\mathbb Z)\).  The lattice kinetic energy is the MMST
nearest-neighbour Hamiltonian, the observable embeddings are induced by the orthonormal
Daubechies \(D4\) mask, and the interaction convention is
\[
 \lambda\eps_N\sum_{x\in\Lambda_N}{:}\Phi_N(x)^4{:}_{c_{N,L}}.
\]
No Daubechies--Galerkin kinetic energy, infinite-line Hamiltonian, or continuum
Hamiltonian is identified with this operator.  Spatial-cutoff and thermodynamic limits
are separate limits and require separately stated hypotheses.
\end{resultbox}

\tableofcontents

\section{Geometric and canonical data}

\LS{1}{1}{Fix the torus and the dyadic lattices.}

Fix \(L>0\), choose \(s_0\in\mathbb N_0\), put \(r_0=2^{s_0}\) and
\[
 \eps_0:=\frac{L}{r_0},\qquad
 \eps_N:=2^{-N}\eps_0,\qquad r_N:=2^Nr_0.                         \tag{1.1}
\]
The position and momentum lattices are
\begin{align}
 \Lambda_N&:=\{\eps_Nn:-r_N\le n\le r_N-1\}\subset\T_L,\\
 \Gamma_N&:=\left\{\frac{\pi j}{L}:-r_N\le j\le r_N-1\right\}.
                                                                    \tag{1.2}
\end{align}
All position indices are taken modulo \(2r_N\).  There are exactly \(2r_N\)
sites and \(2r_N\eps_N=2L\).

\LS{1}{2}{Fix the Fourier normalization.}

For \(f:\Lambda_N\to\C\), define
\[
 (F_Nf)(k):=\sum_{x\in\Lambda_N}f(x)e^{-ikx},\qquad k\in\Gamma_N. \tag{1.3}
\]
Then
\[
 f(x)=\frac1{2r_N}\sum_{k\in\Gamma_N}(F_Nf)(k)e^{ikx},\qquad
 \sum_x|f(x)|^2=\frac1{2r_N}\sum_k|(F_Nf)(k)|^2.                 \tag{1.4}
\]
For the real canonical coordinates below, set
\[
 \widehat q_N:=\eps_N^{1/2}F_Nq,\qquad
 \widehat p_N:=\eps_N^{1/2}F_Np.                                \tag{1.5}
\]

\LS{1}{3}{Fix the symplectic and Weyl conventions.}

Let \(h_{N,L}=\ell^2(\Lambda_N)\).  Write every one-particle vector uniquely as
\[
 \xi=\eps_Nq+ip,\qquad q,p:\Lambda_N\to\R.                     \tag{1.6}
\]
The symplectic form is
\[
 \sigma_{N,L}((q,p),(q',p'))
 =\eps_N\sum_{x\in\Lambda_N}\bigl(q(x)p'(x)-p(x)q'(x)\bigr).    \tag{1.7}
\]
The Weyl algebra \(\W_{N,L}\) is generated by unitaries \(W_N(\xi)\) satisfying
\[
 W_N(\xi)W_N(\eta)=e^{-i\sigma_{N,L}(\xi,\eta)/2}W_N(\xi+\eta). \tag{1.8}
\]
In the regular Schrödinger representation,
\begin{align}
 W_N(\eps_Nq+ip)&=e^{i(\Phi_N(q)+\Pi_N(p))},\\
 \Phi_N(q)&=\eps_N\sum_xq(x)\Phi_N(x),\qquad
 \Pi_N(p)=\eps_N\sum_xp(x)\Pi_N(x),                            \tag{1.9}
\end{align}
and therefore
\[
 [\Phi_N(x),\Pi_N(y)]=i\eps_N^{-1}\delta_{xy}I.                \tag{1.10}
\]

\section{The D4 observable embeddings}

\LS{2}{1}{Fix the low-pass mask, including its normalization.}

Put
\begin{align}
 h_0&=\frac{1+\sqrt3}{4\sqrt2},&
 h_1&=\frac{3+\sqrt3}{4\sqrt2},\\
 h_2&=\frac{3-\sqrt3}{4\sqrt2},&
 h_3&=\frac{1-\sqrt3}{4\sqrt2},                               \tag{2.1}
\end{align}
and \(h_n=0\) otherwise.  Define
\[
 m_0(t):=2^{-1/2}\sum_{n=0}^3h_ne^{-i n t}.                     \tag{2.2}
\]
With \(z=e^{-it}\), direct substitution gives
\begin{align}
 m_0(t)
 &=\frac18\bigl((1+\sqrt3)+(3+\sqrt3)z+(3-\sqrt3)z^2
                 +(1-\sqrt3)z^3\bigr)\\
 &=\left(\frac{1+z}{2}\right)^2
   \underbrace{\frac{(1+\sqrt3)+(1-\sqrt3)z}{2}}_{=:D(t)}.     \tag{2.3}
\end{align}
In particular \(m_0(0)=1\).  Moreover
\[
 |D(t)|^2=2-\cos t,\qquad
 |m_0(t)|^2=\cos^4(t/2)(2-\cos t).                              \tag{2.4}
\]
Writing \(c=\cos t\), expand
\[
 |m_0(t)|^2+|m_0(t+\pi)|^2
 =\frac{(1+c)^2(2-c)+(1-c)^2(2+c)}4=1.                          \tag{2.5}
\]
Indeed, the two cubic polynomials are $2+3c-c^3$ and
$2-3c+c^3$, whose sum is $4$.
Thus \(0\le |m_0(t)|\le1\), and (2.5) is the quadrature-mirror
identity used below.

\LS{2}{2}{Define the one-step refinement map.}

For \(q,p\in\ell^2(\Lambda_N;\R)\), define
\[
 (R_N^{N+1}(q,p))(y)
 :=\sqrt2\sum_{x\in\Lambda_N}\sum_{n=0}^3
       h_n(q(x),p(x))\,\delta_{x+n\eps_{N+1}}(y).               \tag{2.6}
\]
For \(M<N\), set
\[
 R_M^N:=R_{N-1}^N\circ\cdots\circ R_M^{M+1}.                   \tag{2.7}
\]
The QMF identity (2.5), equivalently
\(
 \sum_nh_nh_{n+2j}=\delta_{j0},
\)
gives by direct expansion
\[
 \sigma_{N+1,L}(R_N^{N+1}\xi,R_N^{N+1}\eta)
 =\sigma_{N,L}(\xi,\eta).                                      \tag{2.8}
\]
Hence the rule
\[
 \alpha_M^N(W_M(\xi)):=W_N(R_M^N\xi)                           \tag{2.9}
\]
defines an injective unital \( * \)-homomorphism and
\[
 \alpha_N^K\circ\alpha_M^N=\alpha_M^K.                         \tag{2.10}
\]

\LS{2}{3}{Record the exact support enlargement.}

If \(q,p\) are supported in \(X\subset\Lambda_M\), every site occurring in
\(R_M^N(q,p)\) has the form
\[
 x+\sum_{j=1}^{N-M}n_j\eps_{M+j},\qquad x\in X,\quad
 n_j\in\{0,1,2,3\}.                                           \tag{2.11}
\]
Consequently, before periodization,
\[
 \supp R_M^N(q,p)
 \subset X+[0,3\eps_M(1-2^{-(N-M)})]
 \subset X+[0,3\eps_M].                                       \tag{2.12}
\]
This is the support enlargement to be used in every local estimate.

\LS{2}{4}{Fix the continuum observable embedding.}

Let \(\phi\) be the compactly supported orthonormal \(D4\) scaling function,
normalized by \(\int\phi=1\) and supported in \([0,3]\).  Let
\[
 \phi_{M,x}(y):=\eps_M^{-1/2}\phi((y-x)/\eps_M),                \tag{2.13}
\]
periodized on \(\T_L\).  Define
\[
 R_M^\infty(q,p):=\eps_M^{1/2}\sum_{x\in\Lambda_M}
      (q(x)+ip(x))\phi_{M,x}.                                  \tag{2.14}
\]
The scaling equation gives
\[
 R_N^\infty R_M^N=R_M^\infty.                                 \tag{2.15}
\]
Thus
\[
 \beta_M(W_M(\xi)):=W_{\rm ct}(R_M^\infty\xi)                 \tag{2.16}
\]
is an injective \( * \)-homomorphism and
\[
 \beta_N\circ\alpha_M^N=\beta_M.                              \tag{2.17}
\]

\section{Nearest-neighbour Hamiltonian and Wick convention}

\LS{3}{1}{Fix the free operator.}

Let \(m>0\).  On
\(
 \Hh_{N,L}=L^2(\R^{2r_N},\dd\bm\varphi)
\)
take the Schrödinger realization of (1.10) and define on Schwartz space
\[
 H_{N,L}^{(0)}
 :=\frac{\eps_N}{2}\sum_{x\in\Lambda_N}
 \left[\Pi_N(x)^2+m^2\Phi_N(x)^2
 +\eps_N^{-2}\bigl(\Phi_N(x+\eps_N)-\Phi_N(x)\bigr)^2\right]. \tag{3.1}
\]
Its one-particle dispersion is exactly
\[
 \gamma_N(k)
 =\sqrt{m^2+4\eps_N^{-2}\sin^2(\eps_Nk/2)},\qquad k\in\Gamma_N.\tag{3.2}
\]
The free vacuum state is
\[
 \omega_{N,L}^{(0)}(W_N(\eps_Nq+ip))
 =\exp\!\left[-\frac1{4(2r_N)}\sum_{k\in\Gamma_N}
 \left(\gamma_N(k)^{-1}|\widehat q_N(k)|^2
       +\gamma_N(k)|\widehat p_N(k)|^2\right)\right].          \tag{3.3}
\]

\LS{3}{2}{Fix the Wick covariance with no hidden additive constant.}

Translation invariance gives
\begin{align}
 c_{N,L}
 &:=\omega_{N,L}^{(0)}(\Phi_N(x)^2)\\
 &=\frac1{2r_N}\sum_{k\in\Gamma_N}\frac1{2\eps_N\gamma_N(k)}.
                                                                    \tag{3.4}
\end{align}
The finite-scale Wick polynomial is, by definition,
\[
 {:}\Phi_N(x)^4{:}_{c_{N,L}}
 :=\Phi_N(x)^4-6c_{N,L}\Phi_N(x)^2+3c_{N,L}^2I.                \tag{3.5}
\]
For comparison, after the thermodynamic limit at fixed \(N\),
\[
 c_N=\frac1{4\pi}\int_{-\pi}^{\pi}
 \frac{\dd\theta}{\sqrt{(m\eps_N)^2+4\sin^2(\theta/2)}}
 =\frac1{2\pi}\log\frac8{m\eps_N}+o(1).                      \tag{3.6}
\]
Equation (3.4), not (3.6), is the convention in the fixed-torus Hamiltonian.

\LS{3}{3}{Fix the interacting Hamiltonian.}

For \(\lambda\ge0\) and a function \(g_N:\Lambda_N\to[0,1]\), define
\begin{align}
 V_{N,L}(g_N)
 &:=\lambda\eps_N\sum_{x\in\Lambda_N}
       g_N(x){:}\Phi_N(x)^4{:}_{c_{N,L}},\\
 K_{N,L}(g_N)&:=H_{N,L}^{(0)}+V_{N,L}(g_N).                    \tag{3.7}
\end{align}
The resolution problem treated in the companion documents uses \(g_N\equiv1\).
The factor \(\lambda\), rather than \(\lambda/4!\), is part of the definition.
The second line of (3.7) denotes the self-adjoint operator obtained from the
closed quadratic form constructed below; it is not an assertion about an
unspecified operator domain.

\begin{proposition}[finite-volume operator and ground state]
For every \(N,L,m>0\), \(\lambda\ge0\), and \(0\le g_N\le1\), the polynomial
Schrödinger expression (3.7) has a lower-bounded closed quadratic-form
realization \(K_{N,L}(g_N)\).  This realization has compact resolvent, and its
lowest eigenvalue is simple.
\end{proposition}

\subsection*{Proof in Lamport form}

\LS{1}{1}{Bound the negative Wick term.}

For every \(u\in\R\) and \(a\ge0\), completing the square gives
\[
 u^4-6au^2+3a^2=(u^2-3a)^2-6a^2\ge-6a^2.                      \tag{3.8}
\]
Thus \(V_{N,L}(g_N)\) is bounded below by a finite constant.

\LS{1}{2}{Establish coercivity of the total potential.}

Write \(d=2r_N\), \(c=c_{N,L}\), and let \(U:\R^d\to\R\) be the complete
multiplication potential in (3.7).  Its gradient contribution is nonnegative,
and its remaining, coordinatewise part is
\[
 \eps_N\sum_x f_x(\varphi_x)
 +3\lambda\eps_Nc^2\sum_xg_N(x),\qquad
 f_x(u):=\lambda g_N(x)u^4+
       \left(\frac{m^2}{2}-6\lambda c g_N(x)\right)u^2.        \tag{3.9}
\]
If \(\lambda=0\), then \(f_x(u)=m^2u^2/2\to\infty\) at every site.  Suppose
\(\lambda>0\).  For each site with \(g_N(x)>0\), the polynomial \(f_x\) has
positive leading coefficient and hence tends to \(+\infty\) as \(|u|\to
\infty\); for each site with \(g_N(x)=0\), it is again \(m^2u^2/2\).
Because the set of sites is finite, for every \(A>0\) there is \(R_A<\infty\)
such that
\[
 \max_x|\varphi_x|\ge R_A
 \quad\Longrightarrow\quad
 \sum_x f_x(\varphi_x)\ge A.                                  \tag{3.10}
\]
Indeed, choose \(R_A\) so that
\(f_x(u)\ge A-\sum_{y\ne x}\inf_{v\in\R}f_y(v)\) whenever
\(|u|\ge R_A\), simultaneously for all \(x\).  Since
\(\|\bm\varphi\|\to\infty\) implies
\(\max_x|\varphi_x|\to\infty\), (3.10) proves
\[
 U(\bm\varphi)\longrightarrow+\infty
 \quad\text{as}\quad\|\bm\varphi\|\longrightarrow\infty.   \tag{3.11}
\]
In the special case \(g_N\equiv1\), one also has the explicit inequality
\[
 \sum_x\varphi_x^4\ge\frac1{2r_N}\left(\sum_x\varphi_x^2\right)^2. \tag{3.12}
\]

\LS{1}{3}{Construct the self-adjoint operator by a closed form.}

Choose \(C>0\) so that \(W:=U+C\ge1\), which is possible by (3.8).
On \(L^2(\R^d)\), define
\begin{align}
 \mathcal D(\mathfrak q_C)
 &:=H^1(\R^d)\cap L^2(\R^d,W\dd\bm\varphi),\\
 \mathfrak q_C[\psi]
 &:=\frac1{2\eps_N}\int_{\R^d}|\nabla\psi|^2\dd\bm\varphi
   +\int_{\R^d}W|\psi|^2\dd\bm\varphi.                       \tag{3.13}
\end{align}
This domain is dense because it contains \(C_c^\infty(\R^d)\).  It is complete
for the norm
\(\|\psi\|_{\mathfrak q_C}^2:=\mathfrak q_C[\psi]\): if
\((\psi_j)\) is Cauchy, it converges in \(H^1\) to some \(\psi\), while
\(W^{1/2}\psi_j\) converges in \(L^2\) to some \(v\).  A subsequence of
\(\psi_j\) converges pointwise almost everywhere to \(\psi\), and the same
subsequence has a further subsequence for which \(W^{1/2}\psi_j\to v\)
pointwise almost everywhere.  Hence \(v=W^{1/2}\psi\), proving
\(\psi\in\mathcal D(\mathfrak q_C)\) and form-norm convergence.  Thus
\(\mathfrak q_C\) is closed.  The representation theorem for closed positive
forms yields a unique positive self-adjoint operator \(K_{N,L}(g_N)+C\);
this defines \(K_{N,L}(g_N)\).  If
\(\psi,\zeta\in C_c^\infty(\R^d)\), integration by parts gives
\[
 \mathfrak q_C(\zeta,\psi)
 =\left\langle\zeta,
 \left(-\frac1{2\eps_N}\Delta+U+C\right)\psi\right\rangle.      \tag{3.14}
\]
Consequently the form operator extends exactly the polynomial differential
expression (3.7) on \(C_c^\infty(\R^d)\).

\LS{1}{4}{Prove compactness of the resolvent.}

Let \((\psi_j)\) be bounded in the form norm.  Given \(B>0\), coercivity
(3.11) supplies \(R_B\) such that \(W(\bm\varphi)\ge B\) for
\(\|\bm\varphi\|\ge R_B\).  Therefore
\[
 \int_{\|\bm\varphi\|\ge R_B}|\psi_j|^2\dd\bm\varphi
 \le B^{-1}\int W|\psi_j|^2\dd\bm\varphi
 \le B^{-1}\sup_j\mathfrak q_C[\psi_j].                       \tag{3.15}
\]
On each ball, \((\psi_j)\) is bounded in \(H^1\), so the Rellich compactness
theorem gives an \(L^2\)-convergent subsequence there.  A diagonal subsequence,
together with the uniform tail bound (3.15) and then \(B\to\infty\), is Cauchy
in \(L^2(\R^d)\).  Hence the form-domain inclusion into \(L^2\) is compact.
The resolvent \((K_{N,L}(g_N)+C)^{-1}\) maps \(L^2\) boundedly into the form
domain by the spectral theorem; composing it with the compact inclusion proves
that this resolvent is compact.

\LS{1}{5}{Prove simplicity of the lowest eigenvalue.}

Compact resolvent and lower boundedness give a lowest eigenvalue.  For \(t>0\),
the Feynman--Kac formula gives the integral kernel
\[
 e^{-tK_{N,L}(g_N)}(\bm x,\bm y)
 =\left(\frac{\eps_N}{2\pi t}\right)^{d/2}
   e^{-\eps_N\|\bm x-\bm y\|^2/(2t)}
   \mathbb E_{\bm x,\bm y}^{\,t}
   \left[\exp\!\left(-\int_0^tU(B_s)\dd s\right)\right].     \tag{3.16}
\]
Here \(\mathbb E_{\bm x,\bm y}^{\,t}\) is expectation for the Brownian
bridge with diffusion coefficient \(\eps_N^{-1}\).  Formula (3.16) follows
first for bounded truncations of the continuous potential from the Trotter
product formula and then for \(U\) by monotone form convergence.  A bridge
path is continuous almost surely, so its image is compact and the integral of
the polynomial \(U\) along it is finite.  Every factor on the right of (3.16)
is therefore strictly positive.  Thus \(e^{-tK_{N,L}(g_N)}\) is positivity
improving.

Choose a ground eigenfunction \(f\).  Since
\(|\nabla|f||\le|\nabla f|\) almost everywhere, the form satisfies
\(\mathfrak q_C[|f|]\le\mathfrak q_C[f]\).  The variational definition of the
lowest eigenvalue forces equality, so \(|f|\) is also a ground eigenfunction.
Applying the positivity-improving semigroup to \(|f|\ne0\) shows that the
normalized ground eigenfunction
\(\psi_0:=|f|/\|f\|\) is strictly positive almost everywhere.

Suppose that the ground eigenspace had dimension at least two.  Since the
operator commutes with complex conjugation, that eigenspace contains a real
ground eigenfunction \(h\ne0\) orthogonal to \(\psi_0\).  The strict positivity
of \(\psi_0\) and \(\langle\psi_0,h\rangle=0\) imply that both \(h_+\) and
\(h_-\) are nonzero.  The strictly positive kernel in (3.16) then gives, almost
everywhere,
\[
 |e^{-tK}h|<e^{-tK}|h|.                                       \tag{3.17}
\]
Pairing (3.17) with \(\psi_0\) gives
\[
 e^{-tE_0}\langle\psi_0,|h|\rangle
 <\langle\psi_0,e^{-tK}|h|\rangle
 =e^{-tE_0}\langle\psi_0,|h|\rangle,
\]
a contradiction.  Hence the lowest eigenvalue is simple.  \QedStep

Denote the normalized ground vector by \(\Omega_{N,L,\lambda,g}\), subtract the ground
energy, and set
\[
 \omega_{N,L,\lambda,g}(A)
 :=\langle\Omega_{N,L,\lambda,g},A\Omega_{N,L,\lambda,g}\rangle. \tag{3.18}
\]

\section{Scope ledger}

\begin{longtable}{P{0.31\textwidth}P{0.17\textwidth}P{0.43\textwidth}}
\toprule
Statement & Status & Exact reason\\
\midrule
\endhead
MMST nearest-neighbour free kinematics & \Pass & Equations (1.1)--(1.10) and
(3.1)--(3.3) fix the one-dimensional specialization.\\
\(D4\) wavelet observable embeddings & \Pass & The mask, QMF identity, maps, and
support enlargement are explicit in (2.1)--(2.17).\\
Finite-torus Wick quartic convention & \Pass & Equations (3.4)--(3.7) fix the
covariance, Wick polynomial, lattice measure, and coupling normalization.\\
Interacting resolution limit & \Fail & Defining the Hamiltonians supplies states but
does not prove convergence of their pullbacks.\\
Infinite-volume/spatial-cutoff limit & \Fail & It is a distinct limit and is not a
consequence of a fixed-torus definition.\\
Daubechies--Galerkin kinetic energy & \Fail & It is not the nearest-neighbour operator
(3.1) and is outside this track.\\
\bottomrule
\end{longtable}

\begin{thebibliography}{9}
\bibitem{MMST}
V.~Morinelli, G.~Morsella, A.~Stottmeister, Y.~Tanimoto,
\emph{Scaling limits of lattice quantum fields by wavelets},
Commun. Math. Phys. \textbf{387} (2021), 1299--1367.
\end{thebibliography}

\end{document}
